\subsection{实验目的}
\begin{itemize}
    \item 基于国产 RISC-V 架构香山仿真环境，掌握 RISC-V GNU Compiler Toolchain、AM 裸机运行时环境和 NEMU 模拟器的基本使用方法，理解程序从编译到运行的完整流程。
    \item 对附件提供的向量求和程序进行实现与运行，确保程序功能正确，能够得到正确的向量求和结果。
    \item 通过对程序执行过程和性能指标的观察，掌握程序性能分析的基本方法，理解不同实现方式对程序执行效率的影响。
    \item 在保证计算结果正确的前提下，对向量求和核心代码进行优化，分析优化前后程序指令流和执行周期的变化。
    \item 结合循环外提、减少过程调用、循环展开等优化手段，理解基础程序优化方法的原理及其对执行效率的提升作用。
    \item 通过多种优化方案的实验比较，寻找当前实验环境下性能更优的实现方式，并形成对实验数据和优化结论的分析能力。通过实验对比不同优化策略的效果，加深对程序优化重要性的认识。
\end{itemize}